% Appendix A: Build & Run Instructions

\section{Appendix A: Build \& Run Instructions}
\label{sec:appendix_build_run}

This appendix provides the compilation, testing, packaging, and execution instructions for the Vendor Engine.

\subsection{System Prerequisites}
The application requires the following development tools installed on the local system:
\begin{itemize}
    \item \textbf{Java Development Kit (JDK)}: Version 17 or higher (the project utilizes Java 17 features, such as switch pattern matching).
    \item \textbf{Apache Maven}: Version 3.8 or higher.
\end{itemize}

\subsection{Standard Build and Test Lifecycle Commands}
All compilation and lifecycle phases are managed via the standard Maven terminal wrapper:
\begin{description}
    \item[Clean Workspaces]: Deletes the target compiler outputs to ensure a fresh build.
    \begin{lstlisting}[language=bash]
mvn clean
    \end{lstlisting}
    \item[Compile Sources]: Compiles all source files in the codebase.
    \begin{lstlisting}[language=bash]
mvn compile
    \end{lstlisting}
    \item[Execute Unit Tests]: Runs the JUnit 5 test suite and generates the code coverage data.
    \begin{lstlisting}[language=bash]
mvn test
    \end{lstlisting}
    \item[Package Executable FAT JAR]: Compiles and packages the application, along with its manifest configurations, into a runnable archive located inside the \code{target/} folder.
    \begin{lstlisting}[language=bash]
mvn package
    \end{lstlisting}
\end{description}

\subsection{Running the Application}
Once packaged, the Vendor Engine can be executed from the root directory using the command line:
\begin{lstlisting}[language=bash]
java -jar target/vendor-engine-1.0.jar
\end{lstlisting}
Running the JAR starts the Swing \code{LoginView} interface. By default, the developer simulation panel is displayed concurrently to allow runtime manipulation.

\subsection{Simulating Network Failover During Live Demos}
To test the offline failover and synchronization features dynamically without disconnecting physical network cables, the \code{NetworkMonitorDaemon} reads the \code{network.flag} file every second. 
The presenter can write values to this file during a live demo to simulate connectivity changes:
\begin{itemize}
    \item \textbf{Simulate Connection Lost}: Write ``down'' to the flag file.
    \begin{lstlisting}[language=bash]
# Windows PowerShell
"down" | Out-File -Encoding ascii network.flag

# Unix / Git Bash
echo down > network.flag
    \end{lstlisting}
    This triggers immediate local serialization, updates the GUI banner to red, and pauses the simulator.
    \item \textbf{Simulate Reconnection}: Write ``up'' to the flag file.
    \begin{lstlisting}[language=bash]
# Windows PowerShell
"up" | Out-File -Encoding ascii network.flag

# Unix / Git Bash
echo up > network.flag
    \end{lstlisting}
    This triggers order synchronization to \code{sync\_payload.json}, hides the warning banner, and resumes order ingestion.
\end{itemize}

% Source: README.md:198-217
% Source: Vendor_Engine_Architecture.md:706-724
% Source: src/main/java/com/festival/vendorengine/io/NetworkMonitorDaemon.java:27-34, 128-138
